\chapter{Anexo A}\renewcommand{\tablename}{Cuadro}\label{an:fdtd-fuentes}
\section{Código, perfiles radiales e informes de ajuste}
El estudio se documentó en el repositorio \href{https://github.com/JavierMencos/Resultados-de-Tesis}{Resultados-de-Tesis}. Para identificar el material utilizado en el análisis se fijaron las versiones enumeradas a continuación.
\begin{enumerate}
\item \href{https://github.com/JavierMencos/Resultados-de-Tesis/blob/0dc64c175678ffb762755169c651f998836e8f8a/Tesis_programa.py}{Programa de simulación y extracción de la atenuación}. El archivo \texttt{Tesis\_programa.py} reunió la definición del medio y del monopolo, el cálculo del flujo radial, la selección del modelo de ajuste y la estimación de los intervalos por bloques.
\item \href{https://github.com/JavierMencos/Resultados-de-Tesis/tree/e96d8776842d98bcb3e97faac9887df156888eac/Tesis_Resultados}{Archivos de las diez configuraciones}. En cada nombre, el número anterior a \texttt{px} indicó la resolución en px/mm. Cinco archivos sin extensión contuvieron perfiles separados por comas, con radio, magnitud del flujo radial y perfil compensado en decibelios. Se conservaron estos perfiles para vacío a 5, concreto a 10, músculo a 10 y 20, y agua salada a 40 px/mm.
\item Para las otras cinco configuraciones, los informes de las figuras documentaron las estimaciones y sus intervalos. Se consultaron los correspondientes a agua salada a \href{https://github.com/JavierMencos/Resultados-de-Tesis/commit/823defa}{5}, \href{https://github.com/JavierMencos/Resultados-de-Tesis/commit/adbe296}{10} y \href{https://github.com/JavierMencos/Resultados-de-Tesis/commit/100e2a0}{20 px/mm}, a \href{https://github.com/JavierMencos/Resultados-de-Tesis/commit/209d7d6}{concreto a 5 px/mm} y a \href{https://github.com/JavierMencos/Resultados-de-Tesis/commit/e96d877}{tejido muscular a 5 px/mm}.
\end{enumerate}

\section{Auditoría de las incertidumbres del ajuste}
Se reprodujo el posprocesamiento de los cinco perfiles numéricos conservados. El coeficiente y ambos límites de cada IC coincidieron con el registro histórico a cuatro decimales. Se verificó además una corrección del remuestreo que incluyó la última posición posible de inicio de un bloque. Con $n$ residuos y longitud $b$, se sortearon los enteros entre 0 y $n-b$, ambos incluidos. El programa original había omitido la última posición al usar un límite superior exclusivo.

Con esta corrección, el extremo superior del IC del vacío pasó de 0.0075 a 0.0074 dB/mm. En los otros cuatro perfiles ambos extremos coincidieron a la precisión registrada. El Cuadro~\ref{tab:fdtd-base} conservó el intervalo histórico, cuya interpretación en vacío quedó limitada por la calidad insuficiente del ajuste. Las cinco configuraciones respaldadas exclusivamente por figuras e informes conservaron sus IC publicados. La auditoría, los códigos de análisis y las tablas se conservaron en la \href{https://github.com/JavierMencos/Resultados-de-Tesis/tree/0720a8d4f872b3a98538f2d91427c173715534f0/Tesis_Resultados/Atenuacion}{versión archivada del análisis de atenuación}.

\section{Registros de la campaña de balance}\label{an:balance-registros}
Se archivaron los resultados en \href{https://github.com/JavierMencos/Resultados-de-Tesis/tree/cfe106687c66daf53a6fc653ad23edda9b855882/Tesis_Resultados/Balance}{Tesis\_Resultados/Balance}. Se utilizaron los identificadores del Cuadro~\ref{tab:balance-identificadores}. Cada nombre completo incluyó el medio, el sufijo de resolución y una marca UTC con precisión de microsegundos. El \href{https://github.com/JavierMencos/Resultados-de-Tesis/blob/cfe106687c66daf53a6fc653ad23edda9b855882/Tesis_Resultados/Balance/tablas/indice_corridas.csv}{índice de corridas} relacionó los identificadores con los nombres completos y sus hashes SHA-256. Las dos pruebas breves de instrumentación se conservaron en la subcarpeta \texttt{pruebas}, separadas de las cinco corridas físicas.

\begin{table}[htbp]\centering\small
\caption{Identificación de las corridas físicas de balance. La fecha UTC correspondió al guardado del archivo.}\label{tab:balance-identificadores}
\begin{tabular}{llrl}\hline
ID & Medio & Resolución (px/mm) & Fecha UTC \\ \hline
B-AS05 & Agua salada & 5 & 2026-11-05 21:57:39 \\
B-AS10 & Agua salada & 10 & 2026-11-06 01:12:39 \\
B-TM05 & Tejido muscular & 5 & 2026-11-06 01:17:20 \\
B-CO05 & Concreto & 5 & 2026-11-06 02:57:12 \\
B-VA05 & Vacío & 5 & 2026-11-06 09:46:53 \\ \hline
\end{tabular}\end{table}

Los JSON conservaron las dos muestras, los límites del anillo en mm, la versión del programa, el sistema, el tiempo de pared y la memoria residente máxima. En esos registros, \texttt{interior} se guardó positivo hacia el volumen, mientras que \texttt{exterior}, \texttt{tapa\_superior} y \texttt{tapa\_inferior} se guardaron positivos hacia fuera. El campo \texttt{residuo\_relativo} empleó como denominador el valor absoluto del flujo de la pared interior. Para los cuadros comparativos de la tesis se reconstruyó la entrada total mediante la Ecuación~\ref{eq:balance-normalizado}. Se distinguieron así ambas normalizaciones.

En la subcarpeta \texttt{tablas} se conservaron un CSV por corrida, el cuadro \texttt{balance\_campana.csv} y la comparación \texttt{balance\_refinamiento.csv}. El primero de esos cuadros reunidos incluyó las regiones físicas, ambos tiempos y los diagnósticos temporales. El segundo comparó únicamente corridas del mismo medio y con límites físicos idénticos. Los signos, las unidades y los denominadores se identificaron en los encabezados y en la guía de reproducción.

\section{Procedimiento de reproducción y entorno}
El archivo \texttt{balance.py} conservó los parámetros electromagnéticos y geométricos del simulador y añadió la integración anular. Se sincronizó el campo magnético antes de evaluar los productos con el campo eléctrico, conforme al procedimiento descrito en la \href{https://meep.readthedocs.io/en/latest/Synchronizing_the_Magnetic_and_Electric_Fields/}{documentación de Meep}. Se aplicó el factor de integración cilíndrico $2\pi r$ explícitamente al muestreo puntual.

Se abrió la carpeta de trabajo en VS Code conectado a Ubuntu y se ejecutaron las tareas de la configuración \texttt{.vscode/tasks.json}. Se fijaron a uno los hilos de OpenMP y OpenBLAS. Cada corrida escribió un JSON nuevo al finalizar. Se ejecutó \texttt{resumir\_balance.py} con la ruta de cada JSON para verificar las identidades de potencia y producir las dos muestras en CSV. Luego se utilizó \texttt{comparar\_balances.py} para obtener los cuadros comunes y la sensibilidad entre mallas. El archivo \texttt{SHA256SUMS} permitió verificar la integridad de los códigos y los registros publicados.

El entorno y los recursos se documentaron en la Sección~\ref{sec:costo-balance}. Se empleó además SciPy 1.17.1 para la auditoría de los ajustes. La RSS máxima correspondió al proceso individual.

\section{Convenciones de identificación y unidades}
El prefijo B identificó el balance, AS el agua salada, VA el vacío, CO el concreto y TM el tejido muscular. Los sufijos 05 y 10 indicaron 5 y 10 px/mm. Estos códigos se usaron como abreviaturas de las corridas del Cuadro~\ref{tab:balance-identificadores}.

Las escalas y conversiones se explicaron en la Sección~\ref{sec:unidades-campana}. Se contrastaron con la \href{https://meep.readthedocs.io/en/latest/Introduction/#units-in-meep}{definición de unidades de Meep} y su \href{https://meep.readthedocs.io/en/latest/Materials/}{formulación de materiales}. La escala de potencia se dedujo a partir de una escala de campo $E_*$ y de la integración cilíndrica tridimensional. Los valores almacenados conservaron la normalización de la fuente numérica.


\section{Archivos de la verificación planar}
Se conservaron el código \texttt{verificar\_onda\_plana.py}, los CSV de campos complejos y flujo, los JSON de parámetros y ajuste, el cuadro resumen y los hashes en \href{https://github.com/JavierMencos/Resultados-de-Tesis/tree/main/Tesis_Resultados/Verificacion_1D}{Tesis\_Resultados/Verificacion\_1D}. La carpeta \texttt{diagnostico\_inicial} contuvo los ocho casos de encendido abrupto y su código; \texttt{suave} contuvo los ocho casos utilizados en el Cuadro~\ref{tab:planar}. El índice relacionó cada VP con su nombre fechado. Se documentó el cambio de activación sin sobrescribir los resultados previos.

\section{Fuentes y reproducción del contraste de materiales}\label{an:materiales-fuentes}
Se conservaron el código \texttt{parametros\_referencia.py}, los CSV y el registro de procedencia en \href{https://github.com/JavierMencos/Resultados-de-Tesis/tree/main/Tesis_Resultados/Materiales/Auditoria_20261107}{Tesis\_Resultados/Materiales/Auditoria\_20261107}. El código se ejecutó en VS Code conectado a Ubuntu y utilizó únicamente la biblioteca estándar de Python. Se registraron la versión de IT'IS, los coeficientes de músculo, las condiciones de agua y los hashes de los datos consultados. Se comprobó el límite de salinidad cero contra el modelo de agua pura y la pasividad de los medios.

El informe de auditoría distinguió las fuentes verificadas de las referencias cuya ficha no pudo confirmarse. Las citas originales no verificadas no se utilizaron para atribuir los parámetros. Los cálculos de contraste se identificaron como evaluaciones analíticas y no se incorporaron al inventario de corridas FDTD.

\section{Trazabilidad de los objetivos y del alcance}
El objetivo general y el primer objetivo específico quedaron respaldados para la atenuación dentro de medios homogéneos por los perfiles radiales, la comparación analítica, el refinamiento y la prueba planar. La penetración a través de interfaces no se simuló y quedó fuera de la evidencia obtenida. Para el segundo objetivo se definieron ventanas, unidades, normalización, diagnósticos y un procedimiento reproducible, cuyo alcance fue interno a la campaña. Para el tercero se obtuvo una representación radial relativa, sin validación predictiva frente a datos independientes. Esta delimitación se conservó al evaluar el cumplimiento de los objetivos originales.
